\setcounter{figure}{6}
\begin{figure}[t!]
\centering
\includegraphics[width=0.86\textwidth]{fig_lanmm_field_p1_twotone.pdf}
\caption{\textbf{LaNMM frequency selection: two-tone TI field delivered directly to the deep
$P_1$ (alpha) loop.} As Fig.~\ref{fig:lanmm_p2} but with the field on $P_1$. The deep loop
shows frequency dominance at smaller fields, so the amplitude range is
$\varepsilon\le1$~mV here against $3$~mV there. These finite-amplitude maps extend
beyond the perturbative regime of the theory. Color: signed change of alpha-band power from the $\varepsilon=0$ baseline,
zero-centered; most of the grid is suppression. \emph{(a,b)} $P_1$ and $P_2$ enhancement vs.\ $\varepsilon$ and $\Delta f$; the
black contour in \emph{(a)} outlines the frequency-dominance set
$|f_{\rm out}-\Delta f|<0.3$~Hz (\S\ref{sec:vocab}), computed from the dominant
output frequency because band power alone cannot separate a driven from an autonomous
rhythm; the $\varepsilon=0$ row is excluded (resolution floor $0.5$~Hz). The set widens to
$1.25$~Hz at $\varepsilon=1$~mV.
\emph{(c,d)} Enhancement vs.\ carrier frequency and $\Delta f$ at $\varepsilon=0.5$~mV.
The ridge near alpha persists across the carrier band and includes the autonomous
rhythm. A separate lock-in comparison away from alpha, including its window and
background controls, is given in Table~\ref{tab:lanmm-carrier}; these panels show
band-power changes.}
\label{fig:lanmm_p1}
\end{figure}
